% =============================================================================
% Bloque 6: Metodología (slides 12–16)
% =============================================================================

\section{Metodología}

% --- Slide 12: Pipeline arquitectónico ---
\begin{frame}{Pipeline: del texto real al clasificador entrenado}
\centering
\begin{tikzpicture}[
    every node/.style={font=\scriptsize},
    box/.style={draw=uaiblue, rounded corners, fill=uaibluelight!15,
                minimum width=2cm, minimum height=0.9cm, text width=2cm, align=center},
    box2/.style={draw=uaiorange, very thick, rounded corners, fill=uaiorange!15,
                 minimum width=2cm, minimum height=0.9cm, text width=2cm, align=center},
    box3/.style={draw=uaigreen, very thick, rounded corners, fill=uaigreen!15,
                 minimum width=2cm, minimum height=0.9cm, text width=2cm, align=center},
    arrow/.style={->, >=Stealth, thick, uaiblue!70}
  ]
  % Fila superior
  \node[box]  (real)  at (0,2)    {\textbf{Datos\\reales}\\\tiny $n$-shot 10/25/50};
  \node[box]  (enc)   at (2.6,2)  {\textbf{Encoder}\\\tiny mpnet-base-v2\\\tiny 768 dims};
  \node[box]  (llm)   at (5.2,2)  {\textbf{LLM}\\\tiny Gemini 3 Flash\\\tiny $3\times$ excedente};
  \node[box]  (cands) at (7.8,2)  {\textbf{Candidatas}\\\tiny embeddings};
  \node[box2] (filt)  at (10.4,2) {\textbf{Filtros}\\\tiny geométricos};

  % Fila inferior
  \node[box2] (weight) at (10.4,0)  {\textbf{Pesos}\\\tiny softw.\ $T=0.5$};
  \node[box3] (clf)    at (5.2,0)   {\textbf{Clasificador}\\\tiny LR / SVC / Ridge};
  \node[box]  (eval)   at (0,0)     {\textbf{Macro F1}\\\tiny test fijo};

  % Flechas
  \draw[arrow] (real)  -- (enc);
  \draw[arrow] (enc)   -- (llm);
  \draw[arrow] (llm)   -- (cands);
  \draw[arrow] (cands) -- (filt);
  \draw[arrow] (filt)  -- (weight);
  \draw[arrow] (weight) -- (clf);
  \draw[arrow] (real.south)  .. controls +(0,-0.6) and +(-0.5,0) .. (clf.west);
  \draw[arrow] (clf)   -- (eval);
\end{tikzpicture}

\vspace{0.5em}
\mensaje{Filtrado y ponderación: dos contribuciones desacopladas del LLM.}
\end{frame}

% --- Slide 13: Protocolo experimental (vuelve, antes de los datasets) ---
\begin{frame}{Protocolo experimental: 8.000+ configuraciones}
\small
\begin{columns}[T]

\begin{column}{0.55\textwidth}
\textbf{Datos y modelos:}
\begin{itemize}\setlength\itemsep{0pt}
  \item \textbf{13 datasets} en 9 dominios.
  \item \textbf{$n$-shot:} 10 / 25 / 50 ejemplos por clase.
  \item \textbf{5 clasificadores:} LR, SVC, Ridge, RF, MLP.
  \item \textbf{10 métodos} de aumentación.
  \item \textbf{5 LLMs} (Gemini, GPT-5, Claude, Kimi, GLM).
  \item \textbf{4 modelos} de embeddings.
  \item \textbf{5 semillas:} 42, 123, 456, 789, 1011.
\end{itemize}
\end{column}

\begin{column}{0.45\textwidth}
\textbf{Protocolo estadístico:}
\begin{itemize}\setlength\itemsep{0pt}
  \item Métrica: \textbf{macro F1}.
  \item $\Delta$F1 vs.\ SMOTE (puntos porcentuales).
  \item \textbf{$t$-pareada} con corrección de Bonferroni.
  \item Bootstrap CI 95\% (10\,000 remuestreos).
  \item $d$ de Cohen como tamaño de efecto.
  \item Win-rate como medida de consistencia.
\end{itemize}

\vspace{0.3em}
\setbeamercolor{block title}{fg=white,bg=uaiorange}
\begin{block}{Volumen}
\centering\bfseries\large 8.000+ configuraciones
\end{block}
\end{column}

\end{columns}
\end{frame}

% --- Slide 14: Conjuntos de datos (tabla) ---
\begin{frame}{Conjuntos de datos: 13 corpus en 9 dominios}
\centering
\scriptsize
\renewcommand{\arraystretch}{0.92}
\begin{tabular}{@{}llccr@{}}
\toprule
\textbf{Conjunto} & \textbf{Dominio} & \textbf{Clases} & \textbf{$n$-shots} & \textbf{Test} \\
\midrule
\rowcolor{uaibluelight!12}\multicolumn{5}{l}{\textit{\color{uaibluedark}Clasificación de texto --- experimento principal}} \\
\dataset{sms\_spam}              & Detección de spam        & 2  & 10, 25, 50 & 1.394 \\
\dataset{hate\_speech\_davidson} & Discurso de odio         & 3  & 10, 25, 50 & 4.948 \\
\dataset{20newsgroups}           & Noticias (4 cat.)        & 4  & 10, 25, 50 & 7.532 \\
\dataset{ag\_news}               & Noticias (4 cat.)        & 4  & 10, 25, 50 & 7.600 \\
\dataset{emotion}                & Emociones                & 6  & 10, 25, 50 & 2.000 \\
\dataset{dbpedia14}              & Clasificación de temas   & 14 & 10, 25, 50 & 5.000 \\
\dataset{20newsgroups\_20class}  & Noticias (20 cat.)       & 20 & 10, 25, 50 & 7.532 \\
\midrule
\rowcolor{uaiorange!12}\multicolumn{5}{l}{\textit{\color{uaibluedark}Adicionales --- escalabilidad por número de clases}} \\
\dataset{trec6}                  & Preguntas                & 6   & 10, 25, 50 & 500 \\
\dataset{banking77}              & Intenciones bancarias    & 77  & 10, 25, 50 & 3.080 \\
\dataset{clinc150}               & Intenciones de diálogo   & 150 & 10, 25, 50 & 4.500 \\
\midrule
\rowcolor{uaigreen!12}\multicolumn{5}{l}{\textit{\color{uaibluedark}Reconocimiento de entidades nombradas (NER)}} \\
\dataset{MultiNERD}              & Entidades generales      & 5 tipos & 10, 25, 50 & 10.000 \\
\dataset{WikiANN}                & Entidades Wikipedia      & 3 tipos & 10, 25, 50 & 10.000 \\
\dataset{Few-NERD}               & Entidades detalladas     & 8 tipos & 10, 25, 50 & 10.000 \\
\bottomrule
\end{tabular}

\vspace{0.2em}
{\scriptsize\itshape $n$-shot estricto: 10 / 25 / 50 ejemplos por clase. Rango de complejidad: de binario (\dataset{sms\_spam}) a 150 clases (\dataset{clinc150}).}
\end{frame}

% --- Slide 15: Por qué estos datasets de clasificación ---
\begin{frame}{¿Por qué estos datasets de clasificación?}
\begin{columns}[T]

\begin{column}{0.50\textwidth}
\setbeamercolor{block title}{fg=white,bg=uaibluedark}
\begin{block}{\small Experimento principal --- 7 datasets}
\scriptsize
Buscan \textbf{cobertura en tres ejes}:
\begin{itemize}\setlength\itemsep{0pt}
  \item \textbf{Granularidad:} 2 $\to$ 20 clases.
  \item \textbf{Dominio:} spam, odio, noticias, emoción, ontología.
  \item \textbf{Estilo:} corto/largo, formal/informal, social/curado.
\end{itemize}
\vspace{0.4em}
\begin{itemize}\setlength\itemsep{0pt}
  \item \dataset{sms\_spam} (2): NLP canónico, frase corta.
  \item \dataset{hate\_speech\_davidson} (3): dominio sensible.
  \item \dataset{20newsgroups}, \dataset{ag\_news} (4): noticias.
  \item \dataset{emotion} (6): el ejemplo recurrente.
  \item \dataset{dbpedia14} (14): ontología Wikipedia.
  \item \dataset{20newsgroups\_20class} (20): máxima granularidad.
\end{itemize}
\end{block}
\end{column}

\begin{column}{0.50\textwidth}
\setbeamercolor{block title}{fg=white,bg=uaiorange}
\begin{block}{\small Escalabilidad --- 3 datasets adicionales}
\scriptsize
Responden la pregunta:\\[2pt]
\alert{¿el filtrado sigue funcionando con muchas clases?}
\vspace{0.4em}
\begin{itemize}\setlength\itemsep{0pt}
  \item \dataset{trec6} (6): preguntas TREC, mismo régimen de clases que el principal pero distinto dominio. Sirve de \textbf{control}.
  \item \dataset{banking77} (77): intenciones bancarias, granularidad fina con clases semánticamente cercanas.
  \item \dataset{clinc150} (150): intenciones de diálogo de asistentes virtuales. \textbf{Caso extremo}.
\end{itemize}
\vspace{0.4em}
Permiten medir \textbf{degradación del filtrado} al aumentar el número de clases.
\end{block}
\end{column}

\end{columns}

\mensaje{Diez datasets, dos preguntas: ¿qué tan bueno es el filtrado? ¿hasta dónde escala?}
\end{frame}

% --- Slide 16: ¿Qué es NER y por qué la usamos? ---
\begin{frame}{¿Qué es NER y por qué la usamos?}

\setbeamercolor{block title}{fg=white,bg=uaibluedark}
\begin{block}{\small NER (\textit{Named Entity Recognition})}
\scriptsize
Tarea de NLP que \textbf{detecta nombres propios dentro del texto y los clasifica por tipo}: personas, lugares, organizaciones, fechas, etc.

\vspace{0.3em}
\textit{Ejemplo:} «Llamé a \colorbox{uaiblue!25}{\strut\,María\,} en \colorbox{uaiorange!25}{\strut\,Madrid\,} el \colorbox{uaigreen!25}{\strut\,2 de marzo\,}.»\hspace{0.5em}{\tiny \textcolor{uaiblue}{$\uparrow$PER}\hspace{1em}\textcolor{uaiorange}{$\uparrow$LOC}\hspace{1em}\textcolor{uaigreen}{$\uparrow$DATE}}
\end{block}

\vspace{0.4em}
\begin{columns}[T]

\begin{column}{0.55\textwidth}
\small
\textbf{¿Por qué la usamos como prueba?}
\begin{itemize}\setlength\itemsep{1pt}\scriptsize
  \item Construimos el filtro para \textbf{clasificación}.
  \item NER es una tarea NLP \textbf{distinta}: en vez de etiquetar la oración entera, se etiquetan palabras dentro.
  \item Si el mismo filtro funciona en NER \textbf{sin cambiarlo} $\Rightarrow$ el filtro \alert{no depende de la tarea}.
\end{itemize}

\vspace{0.2em}
{\scriptsize Es la prueba más fuerte de que el método \textbf{generaliza}.}
\end{column}

\begin{column}{0.45\textwidth}
\setbeamercolor{block title}{fg=white,bg=uaigreen}
\begin{block}{\small 3 corpus, distinta complejidad}
\scriptsize
\begin{itemize}\setlength\itemsep{1pt}
  \item \dataset{WikiANN} --- 3 tipos (PER, ORG, LOC). \textbf{Fácil}.
  \item \dataset{MultiNERD} --- 5 tipos (+ DATE, MISC). Medio.
  \item \dataset{Few-NERD} --- 8 tipos jerárquicos. \textbf{Difícil}.
\end{itemize}
\end{block}
\end{column}

\end{columns}

\mensaje{Mismo filtro, otra tarea: si funciona acá, sirve en general.}
\end{frame}

% --- Slide 17: Los 5 filtros ---
\begin{frame}{Cinco estrategias de filtrado evaluadas}
\small
\begin{tabular}{@{}p{2.6cm}p{4.6cm}p{4.6cm}@{}}
\toprule
\textbf{Filtro} & \textbf{Criterio} & \textbf{Parámetros} \\
\midrule
\textit{None} (control)        & Sin filtrado, todas las muestras LLM. & --- \\[2pt]
\textit{LOF}                   & Densidad local (Local Outlier Factor). & $k=20$, threshold $=0$ \\[2pt]
\textit{Cascade L1} \alert{$\star$} & Distancia euclidiana al ancla. Top-$N$ por ranking. & $k=10$, niveles 1--4 \\[2pt]
\textit{Combined}              & LOF \textbf{Y} similitud coseno. & $\tau_{\cos}=0.5$ \\[2pt]
\textit{Embedding-Guided}      & Cobertura $0.6$ + calidad $0.4$, voraz. & multi-criterio \\
\bottomrule
\end{tabular}

\vspace{0.6em}
{\scriptsize\itshape\alert{$\star$} Cascade nivel 1 = mejor configuración encontrada.}

\vfill
\mensaje{Cinco diseños, una pregunta: ¿filtros simples o complejos?}
\end{frame}

% --- Slide 14: Cascade — pieza central ---
\begin{frame}{Cascada de filtros: niveles 1--4}
\begin{columns}[T]

\begin{column}{0.55\textwidth}
\small
\textbf{Score compuesto} (media geométrica):
\[
s(x) = \left(\prod_{i=1}^{L} s_i(x)\right)^{1/L}
\]
\begin{description}\scriptsize
  \item[Nivel 1 \alert{$\star$}] Distancia euclidiana al ancla.
  \item[Nivel 2] Similitud coseno.
  \item[Nivel 3] Pureza KNN ($k$=10).
  \item[Nivel 4] Distancia al centroide.
\end{description}

\vspace{0.3em}
\begin{insight}
\scriptsize Con vectores L2-normalizados, niveles 1 y 2 son redundantes; los niveles 3--4 introducen ruido en pocos ejemplos.
\end{insight}
\end{column}

\begin{column}{0.45\textwidth}
\centering
\begin{tikzpicture}[every node/.style={font=\tiny}]
  % Pirámide de niveles
  \fill[uaiorange!40] (-2.5,2) rectangle (2.5,2.7);
  \node[font=\scriptsize\bfseries] at (0,2.35) {Nivel 1: Distancia $\star$};

  \fill[uaibluelight!30] (-2.2,1.2) rectangle (2.2,1.9);
  \node[font=\scriptsize] at (0,1.55) {Nivel 2: Similitud};

  \fill[uaibluelight!20] (-1.9,0.4) rectangle (1.9,1.1);
  \node[font=\scriptsize] at (0,0.75) {Nivel 3: KNN-purity};

  \fill[uaibluelight!10] (-1.6,-0.4) rectangle (1.6,0.3);
  \node[font=\scriptsize] at (0,-0.05) {Nivel 4: Centroide};

  % Flecha de score
  \draw[->, very thick, uaiorange] (-3.2,-0.4) -- (-3.2,2.7);
  \node[font=\tiny, rotate=90] at (-3.6,1.15) {score $\uparrow$};
\end{tikzpicture}\\[4pt]
{\scriptsize Selecciona top-$N$ por ranking,\\ \textbf{sin umbral binario}.}
\end{column}

\end{columns}
\end{frame}

% --- Slide 15: Ponderación suave ---
\begin{frame}{Ponderación suave: la 2.\textsuperscript{a} contribución}
\begin{columns}[T]

\begin{column}{0.5\textwidth}
\small
En vez de \textbf{aceptar / rechazar}, asignamos un \textit{peso} a cada muestra.

\vspace{0.4em}
\textbf{Eq. 1 — Normalización min-max:}
\[
\tilde{s}(x) = \frac{s(x) - s_{\min}}{s_{\max} - s_{\min}}
\]
\textbf{Eq. 2 — Escalado por temperatura:}
\[
w(x) = \max\!\left(w_{\min},\ \tilde{s}(x)^{\,1/T}\right)
\]

{\scriptsize $T=0.5$, $w_{\min}=0$, datos reales $\Rightarrow$ peso 1.}
\end{column}

\begin{column}{0.5\textwidth}
\centering
\begin{tikzpicture}[every node/.style={font=\tiny}]
  % Eje
  \draw[->,uaigray] (0,0) -- (4.5,0) node[right]{$\tilde{s}(x)$};
  \draw[->,uaigray] (0,0) -- (0,3.2) node[above]{$w(x)$};
  \node[font=\tiny] at (4,-0.25) {1};
  \node[font=\tiny] at (-0.2,3) {1};

  % Curva T=0.5 (peaked) — w = s^2
  \draw[very thick, uaiorange, domain=0:4, samples=40, smooth] plot (\x, {3*(\x/4)^2});
  \node[uaiorange, font=\scriptsize] at (3.6,2.6) {$T=0.5$};

  % Curva T=1 (linear)
  \draw[thick, uaiblue, dashed, domain=0:4, samples=40] plot (\x, {3*(\x/4)});
  \node[uaiblue, font=\scriptsize] at (3.6,1.5) {$T=1$};

  \node[align=center, font=\tiny] at (2.2,-0.7) {Score normalizado};
\end{tikzpicture}\\[4pt]
{\scriptsize $T<1$ amplifica diferencias entre buenas y malas.}
\end{column}

\end{columns}

\mensaje{Las muestras se entregan al clasificador como \texttt{sample\_weight}.}
\end{frame}

